\documentclass[12pt,a4paper]{article}

\usepackage[T1]{fontenc}
\usepackage[utf8]{inputenc}
\usepackage{amsmath,amssymb,amsfonts}
\usepackage{graphicx}
\usepackage{hyperref}
\usepackage{xcolor}
\usepackage{booktabs}
\usepackage{enumitem}
\usepackage{float}
\usepackage{geometry}

\geometry{
  left=3.2cm,
  right=3.2cm,
  top=2.8cm,
  bottom=2.8cm
}

\hypersetup{
  colorlinks=true,
  linkcolor=blue!60!black,
  citecolor=blue!60!black,
  urlcolor=blue!60!black
}

\setlength{\parindent}{0pt}
\setlength{\parskip}{0.6em}

\pagestyle{plain}

\begin{document}

{\small\scshape paper}

\vspace{0.8cm}

{\LARGE Personality Is Not Who You Are.\\[2pt]
It Is Your Cost Tensor.\\[6pt]
\normalsize A four-dimensional, context-dependent reformulation of personality}

\vspace{0.7cm}

Hakvin Vosteen

\vspace{0.6cm}

\noindent\rule{\textwidth}{0.4pt}

\textbf{Abstract}

\vspace{0.2em}

Standard personality psychology assigns each person a small number of scalar trait scores assumed stable across contexts. We replace this object with a four-dimensional cost tensor $\mathbf{c}_i \in \mathbb{R}^4_{\geq 0}$ over the costs of being publicly wrong, ignored, categorised, and rude, and show that observed action follows the simple rule $\text{speak} \Leftrightarrow B_i > \|\mathbf{c}_i\|_1$. The cost tensor depends explicitly on social context through a multiplicative operator $\mathcal{L}_W$ acting on $\mathbb{R}^4$. Within this formulation, the Big Five emerge as a five-dimensional projection $\Pi_5(\mathbf{c}_i)$ that loses both the context dependence and the action mechanism. Trauma corresponds to a single high-magnitude diagonal operator $\Lambda_{\mathrm{trauma}}$ with at least one eigenvalue $\lambda_k \gg 1$ that permanently inflates a cost component; therapy corresponds to an iterated sequence of contractive inverse operators with eigenvalues less than one. We state a structural identification result: personality change requires an external operator sequence acting on the dominant component, and pure self-decision is insufficient. The framework reproduces every Big Five regularity as a special case while supplying the missing context dependence and the missing intervention mechanism.

\vspace{0.4em}

{\small\textbf{Keywords:} personality psychology, decision theory, signaling costs, operators on cost vectors, trauma, behavioural change, cross-domain modelling}

\noindent\rule{\textwidth}{0.4pt}

\vspace{0.5cm}

% ============================================================
\section{Introduction}
\label{sec:intro}
% ============================================================

You speak when expected benefit exceeds the total cost of speaking. That is the only equation that matters:
\begin{equation}
  \text{speak} \quad \Longleftrightarrow \quad B_i > \|\mathbf{c}_i\|_1
  \label{eq:action}
\end{equation}
where $B_i$ is the expected benefit of the action and $\|\mathbf{c}_i\|_1$ is the total cost.

The cost is not a single number. It is a vector
\begin{equation}
  \mathbf{c}_i = \begin{pmatrix} c_{\mathrm{wrong}} \\ c_{\mathrm{reject}} \\ c_{\mathrm{class}} \\ c_{\mathrm{rude}} \end{pmatrix}
  \quad \in \mathbb{R}^4_{\geq 0}
  \label{eq:tensor}
\end{equation}
with components: cost of being publicly wrong, cost of being ignored or mocked, cost of being categorised, and cost of seeming annoying.

The standard personality view assigns each person a small set of scalar trait scores, fixed across contexts. We claim this is the wrong object. The right object is the four-dimensional vector in \eqref{eq:tensor}, evaluated context by context. Introversion is not low energy. Introversion is high $c_{\mathrm{reject}}$ in unfamiliar environments. The same person, the same brain, but a different operator and therefore a different cost tensor.

% ============================================================
\section{The cost tensor under social context}
\label{sec:context}
% ============================================================

Let $\mathcal{L}_W : \mathbb{R}^4 \to \mathbb{R}^4$ denote the operator induced by social environment $W$. Most environments act mildly on $\mathbf{c}_i$, contracting or expanding individual components by small factors. In familiar environments, all four components decrease toward a baseline. In unfamiliar environments, $c_{\mathrm{reject}}$ and $c_{\mathrm{class}}$ rise sharply while $c_{\mathrm{wrong}}$ and $c_{\mathrm{rude}}$ are largely unchanged.

Figure~\ref{fig:cost_context} shows the resulting profiles. The high-$c_{\mathrm{reject}}$ profile produces a person who is silent with strangers and unlocks at the acquaintance threshold. The low-$\|\mathbf{c}_i\|_1$ profile produces a person who speaks at every social distance. Both are rational. Both follow \eqref{eq:action}. The difference is the cost tensor, not the underlying decision rule.

\begin{figure}[H]
  \centering
  \includegraphics[width=0.85\textwidth]{figures/fig1_cost_vs_context.pdf}
  \caption{Total cost $\|\mathbf{c}_i\|_1$ as a function of social familiarity for two profiles. The introvert profile has high $c_{\mathrm{reject}}$ that decays slowly with familiarity. The extrovert profile has uniformly low $\|\mathbf{c}_i\|_1$. The shaded region indicates contexts in which the introvert is silent because cost exceeds the speak threshold; the same person becomes vocal once familiarity crosses the threshold.}
  \label{fig:cost_context}
\end{figure}

% ============================================================
\section{Why the Big Five is not enough}
\label{sec:bigfive}
% ============================================================

\subsection{What the Big Five gets right}

The Big Five \cite{costa1992neo} identifies five stable trait dimensions: Openness, Conscientiousness, Extraversion, Agreeableness, Neuroticism. The framework is validated across more than fifty years of psychometric data and is predictive for career outcomes, health, and relationship quality. It is the best descriptive psychometric framework available.

\subsection{What the Big Five cannot do}

The Big Five assigns one score per trait, independent of context. Introversion score: 0.7. Always. In every room. With every person. But this is wrong, and the wrongness is observable in oneself:

\begin{itemize}[leftmargin=2em,itemsep=0.1em]
\item Loud with close friends, silent with strangers.
\item Confident at work, anxious at parties.
\item Speaks in small groups, freezes in large ones.
\end{itemize}

The Big Five literature labels this variation as measurement noise. The cost tensor framework labels it as the actual signal.

\subsection{The two formalisms compared}

\begin{table}[H]
\centering
\begin{tabular}{lll}
\toprule
                & Big Five                          & Cost Tensor                              \\ \midrule
Unit            & scalar score $\in [0,1]$          & vector $\mathbf{c}_i \in \mathbb{R}^4$  \\
Context         & context-independent               & context-dependent                        \\
Mechanism       & none (descriptive only)           & act $\Leftrightarrow B_i > \|\mathbf{c}_i\|_1$ \\
Intervention    & not specified                     & $\mathcal{L}_W$: operator on $\mathbf{c}_i$ \\
Trauma          & trait shift (unexplained)         & $\Lambda_{\mathrm{trauma}}$: $\lambda_k \gg 1$ \\
Change          & "relatively stable"             & operator sequence required               \\ \bottomrule
\end{tabular}
\caption{The Big Five describes; the Cost Tensor explains. The two are not in opposition: the Big Five is a projection of the Cost Tensor, derived in Section~\ref{sec:remap}.}
\label{tab:compare}
\end{table}

\subsection{Big Five traits as projections of the cost tensor}
\label{sec:remap}

The Big Five traits emerge as special cases of the cost tensor:
\begin{align}
  \text{Neuroticism}      &\equiv \|\mathbf{c}_i\|_1 \text{ globally high} \\
  \text{Agreeableness}    &\equiv c_{\mathrm{rude}} \text{ dominant} \\
  \text{Conscientiousness} &\equiv c_{\mathrm{wrong}} \text{ dominant} \\
  \text{Openness}         &\equiv \|\mathbf{c}_i\|_1 \text{ globally low} \\
  \text{Introversion}     &\equiv c_{\mathrm{reject}} + c_{\mathrm{class}} \text{ dominant}
\end{align}
Formally,
\begin{equation}
  \text{Big Five} = \Pi_5(\mathbf{c}_i),
  \label{eq:projection}
\end{equation}
a projection from the four-dimensional context-dependent object onto five context-independent axes. The projection loses two things: the context dependence (because the projection is averaged across environments) and the action mechanism (because a scalar trait score does not enter \eqref{eq:action}). The Cost Tensor recovers both.

% ============================================================
\section{Trauma is a Whale operator. Therapy is its inverse}
\label{sec:trauma}
% ============================================================

\subsection{How cost components get permanently elevated}

Every social environment applies a Whale Operator $\mathcal{L}_W$ to the cost tensor, contracting or expanding individual components. Most operators are mild and reversible. Trauma is a single high-magnitude operator that permanently inflates one component:
\begin{equation}
  \mathbf{c}_i^{\,\mathrm{after}} = \Lambda_{\mathrm{trauma}} \, \mathbf{c}_i^{\,\mathrm{before}},
  \qquad
  \Lambda_{\mathrm{trauma}} = \mathrm{diag}(\lambda_1, \lambda_2, \lambda_3, \lambda_4),
  \qquad
  \lambda_k \gg 1 \text{ for some } k.
  \label{eq:trauma}
\end{equation}

\subsection{The classroom example}

A teacher mocks a student's answer in front of the class. Single event. The operator scales $c_{\mathrm{reject}}$ by approximately five and $c_{\mathrm{wrong}}$ by approximately three. The student stops raising their hand, not because they have stopped knowing the answers but because $\|\mathbf{c}_i\|_1 > B_i$ from that moment forward. The action threshold has been crossed by the cost side, not the benefit side.

\begin{figure}[H]
  \centering
  \includegraphics[width=0.85\textwidth]{figures/fig2_trauma_therapy.pdf}
  \caption{Time evolution of $\|\mathbf{c}_i\|_1$ across a trauma event followed by therapeutic intervention. Pre-trauma the cost fluctuates near a low baseline below the action threshold. The trauma event applies $\Lambda_{\mathrm{trauma}}$ and pushes the total cost above threshold, where it remains until therapy begins. The dashed segment shows the iterated application of $\mathcal{L}_W^{-1}$ contracting the inflated component back toward baseline.}
  \label{fig:trauma}
\end{figure}

\subsection{Therapy as inverse operator}

Therapy applies $\mathcal{L}_W^{-1}$ iteratively: repeated low-risk exposures that contract the inflated component back toward baseline:
\begin{equation}
  \mathbf{c}_i^{\,\mathrm{healed}} =
  \mathcal{L}_{\mathrm{therapy}}^{(n)} \circ \cdots \circ \mathcal{L}_{\mathrm{therapy}}^{(1)} \, \mathbf{c}_i^{\,\mathrm{trauma}}
  \quad \xrightarrow{\,n \text{ sessions}\,} \quad
  \mathbf{c}_i^{\,\mathrm{baseline}}.
  \label{eq:therapy}
\end{equation}
Each $\mathcal{L}_{\mathrm{therapy}}^{(j)}$ is a contraction with all eigenvalues less than one in magnitude.

\subsection{Identification result}

\textbf{Proposition 1.} \textit{Personality change is possible if and only if a sufficiently strong sequence of Whale operators is applied to the dominant cost component.}

\vspace{0.4em}

The proposition makes the colloquial claim "you cannot just decide to change" mathematically precise. A scalar self-decision is not an operator on $\mathbb{R}^4$. To change $\mathbf{c}_i$ requires either (a) a sustained external sequence (therapy, exposure, environment shift) or (b) a single high-magnitude reversed-direction event with $\lambda_k$ small. Neither has the form of unaided willpower. The framework therefore predicts that interventions targeting the dominant component will succeed and interventions consisting of generic motivation will fail. This is consistent with the established meta-analytic literature on cognitive behavioural therapy versus generic counselling \cite{butler2006empirical}.

% ============================================================
\section{Boundary conditions and prior work}
\label{sec:prior}
% ============================================================

The framework synthesises three streams. Becker \cite{becker1968crime} modelled cost-of-action decisions formally but did not specify a vector-valued cost. Spence \cite{spence1973job} formalised signaling costs in a single dimension. The Big Five literature \cite{costa1992neo,goldberg1993structure} described personality in terms of static averages without context dependence. Edmondson \cite{edmondson1999psychological} measured the outcome (psychological safety in teams) without writing the underlying equation. The new objects in this paper are the Cost Tensor $\mathbf{c}_i \in \mathbb{R}^4$, the Whale Operator $\mathcal{L}_W : \mathbb{R}^4 \to \mathbb{R}^4$, and the Activation Manifold $\mathcal{M}_{\mathrm{act}} = \{(\mathbf{c}_i, B_i) : B_i = \|\mathbf{c}_i\|_1\}$, the surface in $\mathbb{R}^5$ on which a person is exactly indifferent between speaking and remaining silent.

The model is a structural reformulation, not an empirical claim. Section~\ref{sec:tests} states what would empirically distinguish it from the Big Five baseline.

% ============================================================
\section{Testable predictions}
\label{sec:tests}
% ============================================================

\textbf{Prediction 1} (Context interaction). Within-subject variance of behaviour should track within-subject variance of $\mathbf{c}_i$ measured by self-report at the level of the situation, with a coefficient of variation strictly above the Big Five baseline. Falsified if context-tagged self-reports of $\mathbf{c}_i$ explain no more variance than context-independent Big Five scores.

\textbf{Prediction 2} (Therapy mechanism). Treatment effect should scale with the eigenvalue ratio $\lambda_k^{\,\mathrm{therapy}} / \lambda_k^{\,\mathrm{trauma}}$ measured pre-post on the targeted component. Falsified if generic counselling without component targeting produces effects equal to component-targeted exposure.

\textbf{Prediction 3} (Big Five projection). Regressing each Big Five score on the four cost components, averaged across contexts, should yield $R^2$ above $0.7$ for at least four of the five traits. Falsified by failure to reproduce the standard Big Five factor structure.

\textbf{Prediction 4} (Trauma asymmetry). Cost components inflated by trauma persist longer than equivalently-inflated components without an identifiable trauma source. Falsified if natural decay rates are independent of inflation history.

% ============================================================
\section{Boundaries}
\label{sec:boundaries}
% ============================================================

The four components in \eqref{eq:tensor} are stipulated, not derived. The selection covers the cost categories most prevalent in the action-decision literature but is open to revision. A future version may benefit from a fifth axis ($c_{\mathrm{shame}}$) or from collapsing $c_{\mathrm{class}}$ into a sub-component of $c_{\mathrm{reject}}$.

The Whale Operator $\mathcal{L}_W$ is treated as exogenous to the agent. In reality the agent influences the social environment through behaviour and feedback, so $\mathcal{L}_W$ is partially endogenous. A fully closed-loop version of the model would require coupling the action rule \eqref{eq:action} back into the operator dynamics.

The model is silent on group-level phenomena. Group dynamics may amplify or suppress individual cost components nonlinearly; the present formulation only handles the individual decision.

% ============================================================
\section{Conclusion}
\label{sec:conclusion}
% ============================================================

Personality is not a small set of scalar scores. It is a four-dimensional context-dependent cost vector and the operator that transforms it. Behaviour follows the indifference rule $B_i > \|\mathbf{c}_i\|_1$. Standard personality dimensions are projections that lose context and mechanism. Trauma is a single high-eigenvalue diagonal operator. Therapy is a contractive operator sequence. The framework is structural rather than predictive in the standard sense, but it generates testable predictions and identifies an intervention class (component-targeted contractions) that should systematically outperform interventions of the form "decide to change."

\vspace{1.2em}


% ============================================================
\begin{thebibliography}{99}

\bibitem{becker1968crime}
G.~S. Becker,
\textit{Crime and punishment: an economic approach},
Journal of Political Economy \textbf{76} (1968) 169--217.

\bibitem{spence1973job}
M.~Spence,
\textit{Job market signaling},
Quarterly Journal of Economics \textbf{87} (1973) 355--374.

\bibitem{costa1992neo}
P.~T. Costa, R.~R. McCrae,
\textit{Revised NEO Personality Inventory and NEO Five-Factor Inventory professional manual},
Psychological Assessment Resources, Odessa FL (1992).

\bibitem{goldberg1993structure}
L.~R. Goldberg,
\textit{The structure of phenotypic personality traits},
American Psychologist \textbf{48} (1993) 26--34.

\bibitem{edmondson1999psychological}
A.~Edmondson,
\textit{Psychological safety and learning behavior in work teams},
Administrative Science Quarterly \textbf{44} (1999) 350--383.

\bibitem{butler2006empirical}
A.~C. Butler, J.~E. Chapman, E.~M. Forman, A.~T. Beck,
\textit{The empirical status of cognitive-behavioral therapy: a review of meta-analyses},
Clinical Psychology Review \textbf{26} (2006) 17--31.

\end{thebibliography}

\end{document}
